\documentclass[main.tex]{subfiles}


\begin{document}

%from pagenumber to up
% \phantomsection 
% \hypertarget{toc}{}

 

 

% номер секции вручную
\setcounter{section}{44
}
\addtocounter{section}{-1} 

\section{Вынужденные колебания в контуре}


Электромагнитные колебания в колебательном контуре называются \emph{вынужденными}, если они возникают под действием периодической \emph{вынуждающей} силы. 

Так, вынужденные колебания совершаются в контуре, подключенном к источнику \emph{синусоидального} напряжения (рис.\,\ref{fig:p168}).

    \begin{figure}[H]
    \centering
\begin{circuitikz}[circuit ee IEC,font=\small,thin,
    line cap=round,line join=round,
>={Stealth[inset=0pt,round,length=3mm, angle'=20]},
vec/.style={>={Stealth[inset=0pt,round,length=3.5mm, angle'=20]}}]

    \ctikzset{bipoles/americaninductor/coils=3}


% 1 pic

% \path (0,0) --++ (2.5,0) --++ (0,1.5) coordinate (cl) --++ (-2.5,0) --++ (0,-1.5);

%L
% \path (cl) ++ (-{48*\pgflinewidth},0) ++ (0,{.5*\pgflinewidth}) coordinate (cc);
% \fill[yellow] (cc)
% arc (0:180:{13.65*\pgflinewidth})
% arc (0:180:{13.65*\pgflinewidth})
% arc (0:180:{13.65*\pgflinewidth})
% -- cycle
% ;

%C
% \path[] (cl) ++ (0,-1.5) ++ (-{79*\pgflinewidth},{30*\pgflinewidth}) --++ (0,{-60*\pgflinewidth}) ++ ({-20\pgflinewidth},0) coordinate (cc);
% \fill[yellow] (cc) rectangle++ ({20*\pgflinewidth},{60*\pgflinewidth});


\draw (0,0) to[*C,l_=$C$,name=c] ++ (2.5,0) --++ (0,1.5) coordinate (cl); 
\draw[line cap=butt] (cl) to[*L,l_=$L$,mirror] ++ (-2.5,0) coordinate (cll);
\draw (cll) to[*short,-o] ++ (0,-.5) (0,0) to[*short,-o] ++ (0,.5);
\draw[semithick,xshift=-.2cm,scale around={.6:(0.2,.75)}] (0,.75) node[left,inner sep=0.2222em] {$U$} sin ++ (.1,.1) cos ++ (.1,-.1) sin ++ (.1,-.1) cos ++ (.1,.1);
% \node[below] at (c.south) {$q_\text{max}$};
% \node[left] at (c.north west) {$+$};
% \node[right] at (c.north east) {$-$};


\end{circuitikz}
\caption{Вынужденные колебания}
\label{fig:p168}
\end{figure}

Пусть в схеме, показанной на рис.\,\ref{fig:p168}, напряжение источника\footnote{На схеме выводы источника~--- это выколотые точки.} меняется по закону: $U=U_0\sin(\omega t)$. Тогда в контуре происходят колебания заряда и тока с \emph{вынужденной} циклической частотой~$\omega$. (Если бы колебания в контуре были свободными, то они совершались бы с \emph{собственной} циклической частотой~$\omega_\text{с}$ контура, равной\footnote{Подстановка формулы Томсона в определение циклической частоты.}: $\omega_\text{с}=\dfrac{1\strut}{\sqrt{LC}\strut}$.) 

На рис.\,\ref{fig:p169} приведен график зависимости амплитуды силы тока в контуре от циклической частоты (\emph{резонансная кривая}). 

\begin{figure}[H]
    \centering

    \begin{tikzpicture}[font=\small,            line cap=round,                             >={Stealth[inset=0pt,round,length=3mm, angle'=20]},vec/.style={>={Stealth[inset=0pt,round,length=3.5mm, angle'=20]}}]


\draw[->] (0,0) --++ (0,4) node[left,black] {$I_\text{max}$};
\draw[->] (0,0) --++ (6,0) node[below,black] {$\omega$};
\node[below left] (O) at (0,0) {$O$};

\def\X{.6}
\def\R{.2}
\draw[red,thick,domain=0.1:5,samples=1000] plot (\x,{1.5/(sqrt(\R+(\X*\x-1/(\X*\x))^2))});
% \draw[red,dashed,thick,domain=pi:4*pi,yshift={\cy-.3 cm},,xscale=.5,samples=1000] plot (\x,{\xmax*cos(\x r)});
\draw[red,thick] (0,0) -- (.1,{{1.5/(sqrt(\R+(\X*.1-1/(\X*.1))^2))}});


%dash
\draw[gray,densely dashed] ({1/(sqrt(\X*\X))},0) node[below,black] {$\omega_\text{с}$} --++ (0,{1.5/sqrt(\R)});




    \end{tikzpicture}
\caption{Резонансная кривая}
\label{fig:p169}
\end{figure}


Как видно из рисунка, амплитуда тока тем больше, чем ближе вынужденная частота~$\omega$ к собственной частоте~$\omega_\text{с}$ контура. При $\omega = \omega_\text{с}$ (возможно, приблизительно) наступает \emph{резонанс}~--- возрастание амплитуды колебаний. При $\omega\to0$ ток в контуре равен нулю (ток малой частоты~--- это практически постоянный ток, который через конденсатор пройти не может); ток также равен нулю при $\omega\to\infty$ (при быстром изменении тока в катушке возникает большая ЭДС самоиндукции, препятствующая увеличению тока).

В случае колебаний в контуре (свободных или вынужденных) в нем протекает так называемый \emph{переменный ток}~--- ток, изменяющийся с течением времени. 















 
\end{document}